\documentclass[11pt]{article}
\usepackage[a4paper,margin=1in]{geometry}
\usepackage[brazilian]{babel}
\usepackage[utf8]{inputenc}
\usepackage{graphicx}
\usepackage{float}
\usepackage{needspace}
\usepackage{booktabs}
\usepackage{caption}
\usepackage[numbers,sort&compress]{natbib}
\usepackage{hyperref}
\usepackage{amsmath}
\usepackage{longtable}

\title{Medidas de Adaptação e Análise de Custo-Benefício:\\Estudo de Caso Brasil (Porto Alegre)}
\author{
Alisson Tallys Geraldo Fiorentin$^{1}$ \and
Henrique Sanchez Franz$^{2}$ \and
Leonardo Laipelt$^{1}$ \and
Gabriel Matte Rios$^{1}$
}
\date{\today}

\begin{document}
\maketitle

\begin{center}
\footnotesize
$^{1}$Universidade Federal do Rio Grande do Sul (UFRGS) \quad
$^{2}$Universidade Federal do Rio Grande (FURG)
\end{center}
\vspace{1em}

\begin{abstract}
Este relatório documenta a extensão do estudo de caso Brasil do Starter Pack do IDB
RiskHub / CLIMADA (UNU-EHS), cobrindo o risco de inundação em Porto Alegre, Rio Grande do
Sul, para três tipos de ativo: Schools, Residential e Companies. A partir do exemplo já
existente para Schools, o exercício calibrou Residential e Companies, parametrizou seis
medidas de adaptação por ativo (dezoito no total, cobrindo os seis mecanismos de medida do
CLIMADA) e rodou uma análise de custo-benefício para o período de 2025 a 2050. Como a
análise de custo-benefício usa um hazard de profundidade por período de retorno,
Residential e Companies recebem uma curva de dano por profundidade da literatura
\citep{huizinga2017}, compatível em unidade física com esse hazard e consistente com a
curva já usada para Schools. Como a amplitude dessa curva é um padrão regional de
literatura, ela é calibrada localmente ancorando-a na perda real observada no evento de
2024, usando o raster de profundidade por período de retorno cuja extensão alagada mais se
aproxima da extensão desse evento como referência física. Os totais resultantes são: risco
climático total de 320,348 milhões de dólares para Schools, 470,267 milhões para
Residential e 828,925 milhões para Companies. Sob essa estimativa de risco, as seis medidas
de Companies continuam economicamente justificadas, enquanto nenhuma das seis medidas de
Residential, no custo atualmente parametrizado, se paga. Rastreamos essa assimetria a duas
causas que se somam: vários dos custos de medida de Residential refletem orçamentos de
programa em escala de cidade, grandes em relação ao seu risco real, e o risco de
Residential, uma vez ancorado na perda observada, é pequeno em relação ao seu valor total
exposto, muito grande. Este é um achado
relevante para a decisão, não uma falha de modelagem.
\end{abstract}

\newpage
\tableofcontents
\newpage

\section{Introdução}

O exercício intersessional do IDB RiskHub pede que cada participante trabalhe com um ou
mais tipos de ativo calibrados do Starter Pack do seu país, selecione três medidas de
adaptação relevantes para o estudo de caso (participantes do Brasil: três medidas
\emph{adicionais}, já que Schools já vinha com três), parametrize essas medidas no Entity
file com custos de implementação e manutenção realistas e um efeito definido sobre o
cálculo de risco, rode uma análise de custo-benefício usando o notebook de exemplo do
Brasil, refine a parametrização até que os resultados pareçam realistas, e prepare figuras
adequadas para apresentação a tomadores de decisão. Os participantes são ainda encorajados
a calibrar quaisquer ativos restantes não calibrados em seu estudo de caso, e a anotar
ideias para melhorar os Starter Packs.

Para o Brasil, a regra de três medidas adicionais se aplica porque o Entity de Schools já
vinha com três medidas (dragagem do rio, reconstrução de diques e um sistema de alerta
antecipado). Residential e Companies não tinham medidas reais pré-existentes, apenas uma
linha de placeholder inerte, então a regra foi aplicada de forma uniforme aos três tipos de
ativo: três medidas de base mais três adicionais, seis no total por ativo, dezoito no
total.

Este relatório cobre quatro frentes de trabalho: a calibração de Residential e Companies,
que não estava disponível anteriormente no pack; a parametrização de seis medidas de
adaptação por ativo, cobrindo os seis mecanismos de medida do CLIMADA; a derivação e a
calibração local da curva de dano por profundidade usada para Residential e Companies na
análise de custo-benefício; e a produção de figuras e tabelas adequadas para tomadores de
decisão. Uma investigação complementar, motivada pelo resultado de que nenhuma das medidas
de Residential se paga, confirmou que isso não é um problema de reprodutibilidade (Seção
\ref{sec:causa-raiz} e as notas de verificação no Apêndice~\ref{app:notas}).

\section{Metodologia}

\subsection{Ambiente}

A análise usa o CLIMADA versão 6.1.0 \citep{climada2026}, fixada deliberadamente: a versão
7, prevista para o final de 2026, reescreve os módulos Measures e CostBenefit usados ao
longo deste exercício. Quem for reproduzir estes resultados deve fixar a mesma versão em vez
de pegar a mais recente, já que as interfaces de Entity file e de custo-benefício das quais
este exercício depende são justamente as que estão previstas para mudar. Os comandos abaixo
criam esse ambiente do zero e terminam com a própria suíte de testes da biblioteca, que deve
reportar \texttt{OK} antes que qualquer notebook seja executado.

\begin{verbatim}
conda create -n climada_env --override-channels -c conda-forge climada=6.1.0 \
    notebook nbconvert -y
conda activate climada_env
python -m unittest climada.engine.test.test_impact   # deve terminar em "OK"
\end{verbatim}

\subsection{Reprodução da linha de base}
\label{sec:baseline}

O notebook original do pack, \texttt{09\_costbenefit.ipynb}, foi executado sem alterações
para confirmar o ambiente. Ele reproduziu exatamente os números de referência
documentados: risco climático total 320,348 milhões de dólares, risco médio anual 19,0853
milhões de dólares, risco residual menos 165,334 milhões de dólares, com três medidas
(dragagem do rio, reconstrução de diques, sistema de alerta antecipado) sobre o ativo
Schools. Como verificação independente do ambiente desta sessão, o mesmo notebook foi
reexecutado uma segunda vez, a partir de um clone limpo do repositório upstream do Starter
Pack, sem nenhuma modificação de qualquer sessão de trabalho. Ele reproduziu os mesmos três
números exatamente, confirmando que a biblioteca CLIMADA e o método de custo-benefício
funcionam corretamente e não são sensíveis a peculiaridades do ambiente local.

Essa mesma verificação no repositório upstream mostrou que o Starter Pack, como publicado,
traz um exemplo completo de custo-benefício apenas para Schools. Residential e Companies
existem no upstream apenas como Entity files não calibrados, sem medidas e sem qualquer
análise de custo-benefício prévia. O trabalho de calibração e custo-benefício para esses
dois ativos, descrito no restante deste relatório, é portanto uma primeira tentativa, não a
reprodução de um resultado já existente e previamente validado.

\subsection{Calibração de Residential e Companies}
\label{sec:calibration}

Schools já vinha calibrado no pack; Residential e Companies não. O método segue o usado em
\texttt{06a\_calibration\_schools.ipynb}: uma busca em grade bidimensional sobre os
parâmetros \texttt{max\_damage} e \texttt{max\_intensity}, escalando uma curva de primeira
estimativa contra o hazard de duração do evento de 2024, seguida por um refinamento fino
unidimensional depois que a grade grosseira identifica a vizinhança do ótimo. O espaço de
busca é definido de forma independente para cada ativo, em vez de reaproveitado de Schools,
porque os três ativos diferem enormemente em escala de exposição e na razão entre valor
exposto e perda observada. Residential, por exemplo, é uma grade que cobre a cidade
inteira, com 69,3 bilhões de dólares de valor exposto contra 212 milhões de dólares de
perda observada, cerca de 0,3\%, então o espaço de busca é ampliado para permitir
\texttt{max\_damage} até 0,01, bem abaixo da faixa que seria adequada para um ativo com uma
razão de dano maior, como Schools. Companies, com uma base de exposição menor e mais
concentrada, é calibrado sobre seu próprio espaço de busca da mesma forma. Os ajustes
resultantes, resumidos na Tabela \ref{tab:calibracao}, reproduzem a perda observada em 2024
para os dois ativos com margem de erro inferior a 1\%.

\begin{table}[H]
\centering
\caption{Resultados de calibração para Residential e Companies. Apenas o
\texttt{max\_damage} é identificado pelos dados; o \texttt{max\_intensity} não é, pelos
motivos expostos a seguir.}
\label{tab:calibracao}
\footnotesize
\begin{tabular}{lrrrr}
\toprule
Ativo & max\_damage (identificado) & Perda modelada (USD) & Perda observada (USD) & Erro \\
\midrule
Residential & 0,028 & 212.265.415 & 212.170.680 & 0,04\% \\
Companies & 0,296 & 442.034.021 & 438.755.400 & 0,75\% \\
\bottomrule
\end{tabular}
\end{table}

A Tabela \ref{tab:calibracao} reporta deliberadamente apenas o \texttt{max\_damage}, ainda
que a busca em grade ajuste dois parâmetros. O motivo é que só o \texttt{max\_damage} é de
fato identificado pelos dados. Essa é uma propriedade conhecida do método, documentada no
próprio notebook \texttt{06a} do pack: calibrar dois parâmetros livres contra um único alvo
agregado deixa um grau de liberdade, de modo que para cada \texttt{max\_intensity} existe
um \texttt{max\_damage} que reproduz a observação. Nossa própria varredura confirma isso
diretamente: mantendo \texttt{max\_intensity} fixo em valores de 0,25 a 79,75 dias,
cobrindo toda a faixa pesquisada, uma busca fina sobre \texttt{max\_damage} ainda alcança
erro abaixo de 3\% para Residential e abaixo de 0,5\% para Companies em todo valor testado.

O mecanismo por trás dessa degenerescência é específico da severidade do evento de 2024. A
duração mediana de inundação nos pontos expostos é de 48 dias para Residential e 29 dias
para Companies, enquanto a curva escalada atinge seu dano máximo em aproximadamente
\texttt{max\_intensity} dividido por três. Consequentemente, a grande maioria do valor
exposto alagado, cerca de 93\% para Residential e 91\% para Companies nos valores que a
busca selecionou, fica na parte plana e saturada da curva, onde a perda modelada não
depende mais da duração. O modelo colapsa efetivamente no produto entre
\texttt{max\_damage} e o valor exposto alagado total, o que fixa \texttt{max\_damage}
analiticamente em 0,0277 para Residential e 0,2942 para Companies, dentro de 1,2\% e 0,4\%
dos valores da Tabela \ref{tab:calibracao}, e deixa \texttt{max\_intensity} quase sem
restrição.

Dois diagnósticos confirmam que o \texttt{max\_intensity} ajustado não carrega informação.
Os dois ativos pousaram em bordas opostas da caixa permitida, Residential em 28,75 dias, o
maior valor que o limite herdado permite, e Companies em 0,25 dias, o primeiro ponto da
grade: a assinatura de um otimizador deslizando por uma função objetivo plana até encontrar
uma restrição. O valor de Companies também é fisicamente implausível por si só, já que uma
curva saturando em 0,25 dias implicaria que uma empresa alagada por duas horas sofre o
mesmo dano que uma alagada pelos 91 dias inteiros que o evento durou. As curvas armazenadas
mantêm esses valores (28,75 e 0,25 dias) para rastreabilidade com os notebooks, mas eles
são reportados aqui como artefato da busca, não como estimativa. Resolver a ambiguidade
adequadamente exigiria o mesmo tipo de julgamento especializado específico do ativo que o
pack aplica a Schools, o que por sua vez exige evidência local sobre como as perdas
residenciais e comerciais escalam com a duração da inundação, evidência que este exercício
não possui. Nada disso afeta os resultados de custo-benefício da Seção
\ref{sec:resultados}, já que a curva baseada em duração não é a usada ali.

Essas curvas calibradas, armazenadas em \texttt{data/entity\_files\_calibrated/}, são
válidas para reconstruir o evento único de 2024, cujo hazard é medido em duração de
inundação (dias). A análise de custo-benefício, por outro lado, roda sobre um hazard de
profundidade por período de retorno (metros); por isso usa uma curva de dano separada,
baseada em profundidade, descrita na Seção \ref{sec:curva-profundidade}, compatível em
unidade física com esse hazard.

\subsection{Medidas de adaptação e mecanismos}

Cada ativo recebeu seis medidas: três que já vinham no pack ou seguem a mesma família de
mecanismo para Residential e Companies, e três medidas adicionais selecionadas para este
exercício, seguindo a instrução do briefing intersessional de selecionar medidas
relevantes para o estudo de caso e parametrizar explicitamente seu efeito no cálculo de
risco. Em vez de variar apenas os parâmetros numéricos de um único mecanismo, as dezoito
medidas foram desenhadas para cobrir todos os seis mecanismos de adaptação que o CLIMADA
oferece, distribuídas entre os três ativos de modo que cada mecanismo seja demonstrado
pelo menos uma vez. Essa escolha é deliberada: um tomador de decisão comparando medidas
que diferem só no custo aprende menos do que um comparando medidas que intervêm no
cálculo de risco de formas genuinamente diferentes.

\begin{table}[H]
\centering
\caption{Mecanismos de medida do CLIMADA demonstrados neste exercício}
\label{tab:mecanismos}
\begin{tabular}{p{5cm}p{8cm}}
\toprule
Mecanismo & Exemplo usado \\
\midrule
Transformação linear do hazard (a, b) & School floor raising, home flood-proofing, portable barriers \\
Corte de frequência do hazard (RP) & Rebuilding levees, city-wide drainage upgrade, commercial levee \\
Transferência de risco (seguro) & Portfolio, residential parametric e business interruption insurance \\
Troca de Exposures (relocação) & Relocação parcial nos três ativos \\
Transformação linear do MDD & Early warning system, community warning, elevated inventory storage \\
Transformação linear do PAA & Flood-resistant building codes, business continuity planning \\
\bottomrule
\end{tabular}
\end{table}

Para o seguro, os limites de retenção (attach) e cobertura (cover) foram derivados da
curva de excedência do próprio modelo (\texttt{Impact.calc\_freq\_curve}, usando o período
de retorno de 10 anos como attach e o de 100 anos como cover), não de valores arbitrários.
O carregamento de custo aplicado foi de 1,5 vezes o prêmio puro, razão pela qual a razão
benefício-custo do seguro fica sempre próxima de 1 dividido por 1,5, aproximadamente 0,66,
nos três ativos. Esse valor reflete o custo de transferência de volatilidade, não uma
expectativa de economia.

Para a relocação, foi construído um arquivo HDF5 de Exposures no qual as localidades de
maior profundidade no período de retorno de 100 anos têm seu valor reduzido a 40\% do
original, correspondendo a uma relocação parcial de 60\% do valor exposto. Uma segunda
variante foi construída para 2050, já escalada pelo crescimento econômico de 2\% ao ano;
sem esse ajuste, o benefício da medida incluiria indevidamente a remoção do próprio
crescimento econômico do cenário futuro.

\subsection{Curva de dano por profundidade para a análise de custo-benefício}
\label{sec:curva-profundidade}

Residential e Companies recebem uma curva de dano baseada em profundidade para a análise
de custo-benefício, compatível com o hazard de profundidade por período de retorno usado
ali. Seu formato vem da curva publicada pelo Joint Research Centre \citep{huizinga2017}: a
Tabela 3-5, referente a uso residencial na América do Sul, é adotada para Residential, e a
Tabela 3-9, referente a uso comercial na América do Sul, é adotada para Companies. Trata-se
da mesma fonte de literatura e do mesmo padrão já usado, embora sem citação explícita, pelo
Entity de Schools que vem no pack, que usa uma curva equivalente baseada em profundidade.

Como a amplitude de uma curva de literatura regional não é, por si só, validada para Porto
Alegre, ela é calibrada localmente contra o único dado local disponível: a perda real
observada no evento de 2024. Não existe um mapa de profundidade para esse evento, apenas
um de duração, então é preciso um proxy físico. Entre os rasters de profundidade por
período de retorno disponíveis, aquele cuja extensão alagada mais se aproxima da extensão
realmente observada em 2024 serve como esse proxy. Comparando a área alagada entre os
rasters (contagem de pixels com profundidade maior que zero, convertida para
km\textsuperscript{2} considerando a latitude de Porto Alegre), a extensão do evento de
2024 (141,2 km\textsuperscript{2}) fica quase exatamente entre o período de retorno de 100
anos (134,3 km\textsuperscript{2}) e o de 500 anos (148,0 km\textsuperscript{2}), mais
próxima ainda do período de retorno de 200 anos (140,9 km\textsuperscript{2}, uma
diferença de apenas 0,3 km\textsuperscript{2}). Isso é consistente com o relato público
sobre as enchentes do Rio Grande do Sul em 2024, um evento excepcional, próximo do recorde
histórico para Porto Alegre, e o raster de 200 anos é usado como âncora.

A amplitude da curva é calibrada por reescala: mantendo os mesmos pontos de quebra de
profundidade e o mesmo formato relativo da curva JRC, cada valor de MDD (mean damage
degree) é multiplicado por um único fator de escala, calculado como a razão entre a perda
observada em 2024 e a perda que a curva sem ajuste produziria sobre o raster de período de
retorno de 200 anos (0,048339 para Residential; 0,360351 para Companies). Como nenhum
valor da curva reescalada satura em MDD igual a 1 dentro da faixa de profundidades
observadas (0 a cerca de 9 metros), essa reescala se propaga de forma exatamente linear
pelo resto da análise: o risco total, o risco médio anual e o benefício de cada medida (com
exceção do seguro) são, cada um, o valor que a curva sem ajuste produziria, multiplicado
pelo mesmo fator de escala do ativo. Os limites de retenção e cobertura do seguro
paramétrico são calculados a partir da curva de excedência dessa curva calibrada, seguindo
o mesmo método da Seção 2.4.

Os totais resultantes são 470,267 milhões de dólares para Residential e 828,925 milhões
para Companies (Seção \ref{sec:resultados}). Os custos de medida são parametrizados de
forma independente, como orçamentos de obra ou de programa municipal, não como uma
quantidade derivada da escala de risco; a Seção \ref{sec:causa-raiz} examina em detalhe a
consequência dessa independência: nenhuma das seis medidas de Residential se paga aos
custos parametrizados, enquanto todas as seis de Companies continuam se pagando. Essa é
uma limitação explícita a ser levada em conta: a curva é ancorada em um único ponto de
dados agregado (a perda total de um evento), não em uma calibração espacialmente
distribuída como a que orientou o hazard de duração na Seção \ref{sec:calibration}, a mesma ressalva que se
aplica a qualquer calibração construída sobre poucos pontos de referência.

Como conferência independente de que esse fator de escala não é uma escolha arbitrária ou
forçada, a mesma calibração foi reproduzida usando o próprio framework de calibração do
CLIMADA (\texttt{climada.util.calibrate}), que minimiza o erro quadrático médio entre
perda modelada e observada sobre um parâmetro de escala livre, usando
\texttt{scipy.optimize}. Três algoritmos de otimização (Nelder-Mead, L-BFGS-B e Powell)
foram rodados de forma independente e todos convergiram para os mesmos valores obtidos
analiticamente, 0,048339 para Residential e 0,360351 para Companies, confirmando que esses
são os valores que minimizam o custo para esse modelo de um parâmetro, dada a perda
observada em 2024, não um valor escolhido para produzir um resultado específico.

\subsection{Sensibilidade ao período de retorno usado como âncora}
\label{sec:sensibilidade}

O fator de escala é sensível a qual período de retorno é usado como proxy físico do
evento de 2024: calculado da mesma forma para cada período de retorno disponível, ele varia
de 2,37 (Residential, RP10) até 0,039 (Residential, RP500), já que a perda modelada pela
curva crua cresce com a profundidade enquanto a perda observada em 2024 é fixa. Isso torna
a escolha do RP200, baseada na semelhança de área alagada (Seção \ref{sec:curva-profundidade}),
a decisão metodológica mais consequente de toda esta análise, e vale a pena conferir o
quanto os resultados relatados mudariam sob uma escolha diferente, ainda assim fisicamente
plausível.

\begin{table}[H]
\centering
\caption{Risco total e B/C das medidas não-seguro sob períodos de retorno alternativos como âncora}
\label{tab:sensibilidade}
\small
\begin{tabular}{lrrr}
\toprule
& RP100 & RP200 (usado) & RP500 \\
\midrule
Risco total Residential (M USD) & 569,0 & 470,3 & 378,1 \\
Faixa de B/C, medidas de Residential & 0,23 a 0,88 & 0,19 a 0,73 & 0,15 a 0,59 \\
Risco total Companies (M USD) & 1.019,8 & 828,9 & 632,5 \\
Faixa de B/C, medidas de Companies (exceto relocação) & 4,37 a 11,17 & 3,55 a 9,08 & 2,71 a 6,93 \\
Relocação de Companies, B/C & 1,62 & 1,32 & 1,01 \\
\bottomrule
\end{tabular}
\end{table}

Em toda essa faixa, que vai de um evento de 100 anos a um de 500 anos, a conclusão
relevante para a decisão é estável: todas as medidas não-seguro de Residential ficam
abaixo de uma razão benefício-custo de 1, e todas as medidas de Companies, exceto a
relocação, ficam acima de 1. A relocação de Companies é o único resultado deste relatório
sensível o suficiente para merecer destaque isolado: sua razão vai de 1,62 em RP100 para
1,01 em RP500, ficando do lado positivo do ponto de equilíbrio em toda a faixa testada, mas
com uma margem estreita e dependente da âncora escolhida. A razão do seguro é excluída
dessa tabela porque depende da curva de excedência do modelo em vez de escalar
linearmente com a âncora escolhida; na prática ela se mantém perto de 0,66 em todo período
de retorno testado, já que é definida pelo carregamento de custo de 1,5 vez, não pelo nível
de risco subjacente.

\subsection{Critérios de aceite}
\label{sec:refinamento}

A parametrização de cada medida é verificada segundo três critérios: nenhuma razão
benefício-custo infinita, indefinida ou superior a 50, o que sinalizaria um custo
irrealisticamente baixo; nenhum benefício de medida individual que exceda o risco
climático total do ativo; e limites de seguro sempre derivados do modelo, nunca de
estimativas arbitrárias. Os dois primeiros são verificações de sanidade contra um custo mal
especificado ou um benefício calculado contra a linha de base errada; o terceiro existe
porque parâmetros de transferência de risco escolhidos manualmente podem produzir quase
qualquer prêmio e, com ele, quase qualquer razão benefício-custo, sem que o erro seja
visível em lugar nenhum da saída.

Na parametrização final, nenhuma razão benefício-custo ultrapassa 9,08, bem abaixo do teto
de 50, e nenhum benefício individual excede o risco climático total do seu ativo; os dois
critérios de sanidade estão satisfeitos. Em Residential, todas as seis razões caem abaixo
de 1,0, o que não é um sinal de erro de parametrização: é a conclusão honesta de que, nessa
estimativa de risco, essas medidas específicas, nos custos em que foram parametrizadas, não
se pagam dentro do horizonte do estudo. Reajustar os custos para forçar razões acima de 1
substituiria a melhor estimativa disponível de custo e risco pelo resultado desejado, que é
exatamente o que esses critérios existem para evitar.

\section{Resultados}
\label{sec:resultados}

Cada uma das três subseções a seguir cobre um tipo de ativo por vez: Schools, Residential
e Companies. Todas seguem a mesma estrutura, para que os resultados entre ativos possam
ser comparados diretamente. Cada uma abre com o risco climático total do ativo e o risco
médio anual para 2050, depois uma tabela listando as seis medidas de adaptação com seu
mecanismo, custo, benefício e razão benefício-custo, depois as cinco figuras geradas para
aquele ativo (o gráfico de custo-benefício, benefícios por período de retorno, o waterfall
de risco evitado, o mapa de risco anual esperado e o mapa de profundidade com pontos de
relocação), e fecha com um parágrafo interpretando o que essas figuras mostram. Uma
subseção final, Síntese Comparativa, extrai os padrões que só ficam visíveis ao comparar
os três ativos lado a lado, como a estabilidade da razão do seguro e a divisão entre clima
e crescimento.

\subsection{Schools}

Schools é o único ativo neste exercício cuja calibração e curva de dano nunca foram
tocadas: seus números de risco abaixo são exatamente os que o Starter Pack já trazia,
sem influência do trabalho de calibração e derivação de curva descrito na Seção 2 para
Residential e Companies. O risco climático total é 320,348 milhões de dólares, com um
risco médio anual de 19,0853 milhões de dólares; os dois números são reproduzidos
exatamente pelo notebook upstream não modificado, como descrito na Seção
\ref{sec:baseline}. A Tabela \ref{tab:schools} lista as seis medidas, três que já vinham
no pack (dragagem do rio, reconstrução de diques e o sistema de alerta antecipado) e três
adicionadas para este exercício (elevação do piso escolar, seguro de portfólio e
relocação parcial), cobrindo sozinhas cinco dos seis mecanismos de medida do CLIMADA.

\begin{table}[H]
\centering
\caption{Schools: medidas de adaptação e resultados de custo-benefício}
\label{tab:schools}
\small
\begin{tabular}{p{3.6cm}p{3cm}rrr}
\toprule
Medida & Mecanismo & Custo (M USD) & Benefício (M USD) & B/C \\
\midrule
River dredging & hazard $-1,5$ m & 789,562 & 195,782 & 0,25 \\
Rebuilding levees & corte RP $\approx$ 100a & 161,171 & 257,866 & 1,60 \\
Early warning system & MDD $\times 0,9$ & 1,500 & 32,035 & 21,36 \\
School floor raising & hazard $-0,4$ m & 7,625 & 64,244 & 8,43 \\
Portfolio flood insurance & seguro 32,3M xs 116,9M & 258,463 & 171,475 & 0,66 \\
Partial relocation (5 de 122) & troca de Exposures & 10,500 & 34,764 & 3,31 \\
\bottomrule
\end{tabular}
\end{table}

Os custos na Tabela \ref{tab:schools} variam de 1,5 milhão de dólares para o sistema de
alerta antecipado até 789,6 milhões para a dragagem do rio, e as razões benefício-custo
variam ainda mais, de 0,25 para a dragagem até 21,36 para o sistema de alerta antecipado,
os dois extremos de um padrão que se repete nos três ativos: medidas de transformação de
hazard sem teto de proteção são caras em relação ao risco que evitam, enquanto medidas
baratas e bem direcionadas alcançam as maiores razões justamente porque sua base de custo
é pequena. As cinco figuras a seguir mostram esse padrão de ângulos diferentes, da
comparação bruta de custo-benefício até como a contribuição de cada medida muda com a
severidade do evento.

\needspace{3.08in}
\begin{figure}[H]
\centering
\includegraphics[width=0.55\textwidth]{figures/brazil_schools_cost_benefit.png}
\caption{Resultados de custo-benefício, Schools}
\label{fig:schools-cb}
\end{figure}

A Figura \ref{fig:schools-cb} representa cada medida como um retângulo cuja largura é o
valor presente líquido do dano evitado e cuja altura é a razão benefício-custo, de modo que
a área não é proporcional a nenhuma das duas isoladamente e as duas dimensões precisam ser
lidas juntas. O sistema de alerta antecipado é a barra alta e estreita à esquerda: uma
razão excelente construída sobre um benefício absoluto muito pequeno. A dragagem do rio é o
oposto, a barra mais larga do gráfico, mas tão baixa que fica abaixo da linha de equilíbrio.
Para um tomador de decisão, esse é o trade-off central do portfólio de Schools, já que a
medida que evita mais dano em termos absolutos é também a que não justifica seu custo.

\needspace{2.73in}
\begin{figure}[H]
\centering
\includegraphics[width=0.55\textwidth]{figures/brazil_schools_benefits_by_return_period.png}
\caption{Benefícios por período de retorno, Schools}
\label{fig:schools-rp}
\end{figure}

A Figura \ref{fig:schools-rp} decompõe os mesmos benefícios por severidade do evento, e
explica por que a dragagem se comporta de forma tão diferente das demais. Sendo uma
transformação de hazard sem teto de proteção, ela continua crescendo em termos absolutos
conforme os eventos ficam mais severos, e domina o benefício no período de retorno de 200
anos. As levees, com corte próximo do período de retorno de 100 anos, e o seguro, entre as
duas, absorvem a maior parte do risco médio no meio da faixa. A relocação e o sistema de
alerta antecipado aparecem apenas como faixas finas ao longo de todo o gráfico, o que é
coerente com seus benefícios absolutos pequenos mesmo onde suas razões são favoráveis.

\needspace{3.29in}
\begin{figure}[H]
\centering
\includegraphics[width=0.48\textwidth]{figures/brazil_schools_waterfall_averted_risk.png}
\caption{Waterfall de risco evitado, Schools}
\label{fig:schools-wf}
\end{figure}

A Figura \ref{fig:schools-wf} separa os motores do risco futuro, em vez do efeito das
medidas. Dos 19 milhões de dólares de risco projetado para 2050, 6 milhões já existiam em
2025, mais 4 milhões vêm do crescimento econômico e mais 10 milhões vêm da mudança
climática. Cerca de 71\% do crescimento do risco entre hoje e 2050 é, portanto, atribuível
ao cenário climático, e não à construção de mais ativos em áreas expostas, uma proporção
que se repete de perto nos outros dois ativos e é discutida na Seção \ref{sec:sintese}.

\needspace{4.26in}
\begin{figure}[H]
\centering
\includegraphics[width=0.55\textwidth]{figures/brazil_schools_map_expected_annual_impact.png}
\caption{Mapa de risco anual esperado, Schools}
\label{fig:schools-eai}
\end{figure}

A Figura \ref{fig:schools-eai} passa do agregado para o espacial, mostrando onde o risco
anual esperado de fato recai. O risco se concentra nos bairros Farrapos e Humaitá e no
entorno do Sarandi, as franjas mais baixas junto ao Lago Guaíba, com o restante da rede
escolar contribuindo muito pouco. Como Schools é um conjunto de 122 instalações pontuais, e
não uma grade, escolas individuais de alto risco aparecem como marcadores discretos, o que
é justamente o que torna viável especificar uma medida direcionada como a relocação.

\needspace{4.99in}
\begin{figure}[H]
\centering
\includegraphics[width=0.55\textwidth]{figures/brazil_map_rp100_depth_schools.png}
\caption{Profundidade no período de retorno de 100 anos e pontos de relocação, Schools}
\label{fig:schools-depth}
\end{figure}

A Figura \ref{fig:schools-depth} sobrepõe a profundidade de inundação de 100 anos às
localizações das escolas e marca as cinco selecionadas para relocação parcial. Elas ficam
exatamente onde a mancha de profundidade é mais escura, que é o critério pelo qual foram
escolhidas, com um caso isolado mais ao sul junto ao lago. Esse agrupamento explica por que
relocar apenas cinco de 122 escolas ainda devolve uma razão benefício-custo acima de três: o
risco nesse ativo é concentrado o suficiente para que uma intervenção pequena e bem
direcionada capture uma fatia desproporcional dele.

\subsection{Residential}

Residential é o ativo mais afetado pela derivação da curva de dano por profundidade na
Seção \ref{sec:curva-profundidade}, já que é calibrado e avaliado pela primeira vez neste
exercício, em vez de reaproveitar um exemplo já existente no pack. O risco climático total
é 470,267 milhões de dólares, com um risco médio anual em 2050 de 28,0255 milhões de
dólares, valores já ancorados na perda real de 2024 em Porto Alegre e conferidos quanto à
sensibilidade à escolha da âncora na Seção \ref{sec:sensibilidade}. A Tabela
\ref{tab:residential} lista as seis medidas parametrizadas para esse ativo: home
flood-proofing e uma modernização de drenagem em escala de cidade, ambas herdando o
mecanismo do exemplo de Schools; uma camada de seguro paramétrico dimensionada para a
curva de excedência própria deste ativo; relocação parcial das residências de maior risco;
e duas medidas adicionadas especificamente para Residential, um sistema de alerta
antecipado comunitário e códigos de obra resistentes a inundação para reformas.

\begin{table}[H]
\centering
\caption{Residential: medidas de adaptação e resultados de custo-benefício}
\label{tab:residential}
\small
\begin{tabular}{p{3.6cm}p{3cm}rrr}
\toprule
Medida & Mecanismo & Custo (M USD) & Benefício (M USD) & B/C \\
\midrule
Home flood-proofing & hazard $-0,3$ m & 362,064 & 70,321 & 0,19 \\
Residential parametric insurance & seguro 47,8M xs 175,4M & 382,446 & 251,631 & 0,66 \\
City-wide drainage upgrade & corte RP $=20$a & 300,000 & 218,109 & 0,73 \\
Relocation (50 de 20.187) & troca de Exposures & 8,250 & 3,924 & 0,48 \\
Community-based early warning & MDD $\times 0,85$ & 150,000 & 70,540 & 0,47 \\
Flood-resistant building codes & PAA $\times 0,8$ & 200,000 & 94,053 & 0,47 \\
\bottomrule
\end{tabular}
\end{table}

Nenhuma dessas seis medidas se paga ao custo parametrizado, nessa estimativa de risco.
Este é o achado central deste relatório, relevante para a decisão, e é examinado em
detalhe na Seção \ref{sec:causa-raiz}, que o rastreia a duas causas que se somam: a escala
do orçamento de várias dessas medidas em relação ao risco real de Residential, e a
diferença subjacente entre quanto valor está exposto e quanto desse valor está realmente
em risco, uma vez que a curva de dano é ancorada numa perda real observada em vez de
deixada na sua amplitude de literatura. As cinco figuras abaixo mostram a mesma conclusão
por ângulos diferentes, da comparação bruta de custo-benefício até como o próprio risco
está distribuído pela cidade.

\needspace{3.10in}
\begin{figure}[H]
\centering
\includegraphics[width=0.55\textwidth]{figures/brazil_residential_cost_benefit.png}
\caption{Resultados de custo-benefício, Residential}
\label{fig:residential-cb}
\end{figure}

A Figura \ref{fig:residential-cb} torna o resultado principal imediatamente visível: todos
os retângulos ficam abaixo da linha de equilíbrio, ou seja, nenhuma medida deste portfólio
devolve mais do que custa. Vale ler o gráfico ao lado do equivalente de Schools (Figura
\ref{fig:schools-cb}), onde as barras cobrem uma ampla faixa de alturas. Aqui elas estão
comprimidas numa faixa estreita entre 0,19 e 0,73, que é a assinatura visual de orçamentos
de medida grandes em relação ao risco que endereçam, e não de uma única medida mal
especificada.

\needspace{2.66in}
\begin{figure}[H]
\centering
\includegraphics[width=0.55\textwidth]{figures/brazil_residential_benefits_by_return_period.png}
\caption{Benefícios por período de retorno, Residential}
\label{fig:residential-rp}
\end{figure}

A Figura \ref{fig:residential-rp} mostra um padrão nítido de corte. A modernização de
drenagem em escala de cidade, com corte no período de retorno de 20 anos, cobre a totalidade
do risco nos períodos de 5 e 10 anos, abaixo do seu teto de proteção, e depois desaparece
completamente a partir do período de 20 anos, quando os eventos excedem esse teto. Daí em
diante, o seguro, na camada de risco mais alta, e os códigos de obra assumem. No período de
retorno de 200 anos sobra uma fatia sem cobertura entre o topo das medidas empilhadas e o
contorno do risco total, mostrando que as seis medidas juntas ainda deixam um resíduo de
risco visível contra o evento mais extremo.

\needspace{3.32in}
\begin{figure}[H]
\centering
\includegraphics[width=0.48\textwidth]{figures/brazil_residential_waterfall_averted_risk.png}
\caption{Waterfall de risco evitado, Residential}
\label{fig:residential-wf}
\end{figure}

A Figura \ref{fig:residential-wf} decompõe o crescimento do risco até 2050. O risco sobe de
8,2 milhões de dólares em 2025 para 28,0 milhões em 2050, um aumento de 5,3 milhões
atribuível ao crescimento econômico e de 14,5 milhões à mudança climática, cerca de 73\%
clima contra 27\% crescimento. O fato de isso coincidir tão de perto com a divisão de
Schools, apesar de uma base de exposição e uma calibração completamente diferentes, é um dos
achados mais robustos do exercício e é retomado na Seção \ref{sec:sintese}.

\needspace{4.20in}
\begin{figure}[H]
\centering
\includegraphics[width=0.55\textwidth]{figures/brazil_residential_map_expected_annual_impact.png}
\caption{Mapa de risco anual esperado, Residential}
\label{fig:residential-eai}
\end{figure}

A Figura \ref{fig:residential-eai} mostra um padrão espacial bem distinto do de Schools. Em
vez de um punhado de pontos concentrados, a exposição residencial cobre quase toda a mancha
urbana de Porto Alegre, porque esse ativo é um dado gradeado de cidade inteira com mais de
20.000 pontos, e não uma lista de instalações. Essa amplitude é justamente o que torna as
medidas em escala de cidade caras aqui: uma intervenção definida sobre toda a grade precisa
ser custeada sobre toda a grade, enquanto o risco que ela endereça permanece concentrado no
norte, na área mais baixa.

\needspace{4.99in}
\begin{figure}[H]
\centering
\includegraphics[width=0.55\textwidth]{figures/brazil_residential_map_rp100_depth_relocation.png}
\caption{Profundidade no período de retorno de 100 anos e pontos de relocação, Residential}
\label{fig:residential-depth}
\end{figure}

A Figura \ref{fig:residential-depth} marca os 50 pontos de grade selecionados para relocação
parcial contra a mancha de profundidade de 100 anos. Diferente de Schools, onde os locais
selecionados formam um agrupamento compacto, estes se espalham por vários bairros
ribeirinhos, com bolsões adicionais mais ao sul. Essa dispersão maior é uma das razões pelas
quais a relocação tem desempenho pior aqui do que nos outros dois ativos: a mesma regra
conservadora, mover 60\% do valor nos pontos mais profundos, captura uma fatia menor do
risco total quando esse risco é difuso em vez de concentrado.

\subsection{Companies}

Companies compartilha o método de calibração e a derivação da curva de dano por
profundidade com Residential, mas difere fortemente em quão concentrada é sua exposição:
uma base comercial menor e mais densa ao longo da orla e no centro histórico, em vez de
uma grade cobrindo a cidade inteira. O risco climático total é 828,925 milhões de
dólares, com um risco médio anual em 2050 de 49,6197 milhões de dólares, ambos já ancorados
na perda observada de 2024 da mesma forma que Residential. A Tabela \ref{tab:companies}
lista as seis medidas parametrizadas para esse ativo: barreiras portáteis contra
inundação e uma levee protegendo o distrito comercial, espelhando os mecanismos de
transformação de hazard e corte de Residential; seguro de interrupção de negócios;
relocação das empresas de maior risco; armazenamento elevado de estoque; e planejamento
de continuidade de negócios e diversificação da cadeia de suprimentos.

\begin{table}[H]
\centering
\caption{Companies: medidas de adaptação e resultados de custo-benefício}
\label{tab:companies}
\small
\begin{tabular}{p{3.6cm}p{3cm}rrr}
\toprule
Medida & Mecanismo & Custo (M USD) & Benefício (M USD) & B/C \\
\midrule
Portable flood barriers & hazard $-0,3$ m & 36,533 & 147,182 & 4,03 \\
Business interruption insurance & seguro 66,7M xs 356,6M & 713,520 & 474,347 & 0,66 \\
Relocation (20 de 4.366) & troca de Exposures & 10,250 & 13,508 & 1,32 \\
Levee (distrito comercial) & corte RP $=20$a & 50,000 & 321,952 & 6,44 \\
Elevated inventory storage & MDD $\times 0,8$ & 18,266 & 165,785 & 9,08 \\
Business continuity and supply chain diversification & PAA $\times 0,7$ & 70,000 & 248,677 & 3,55 \\
\bottomrule
\end{tabular}
\end{table}

Diferente de Residential, todas as seis medidas se pagam (B/C $>1$), consequência direta
da base de exposição de Companies ser muito menor e mais concentrada em relação ao seu
risco (a Seção \ref{sec:causa-raiz} quantifica esse contraste diretamente). A relocação,
em 1,32, se paga com a margem mais apertada das seis, e é o único resultado deste
relatório sensível o suficiente ao período de retorno usado como âncora (Seção
\ref{sec:sensibilidade}) para merecer destaque como menos certo que o restante; as outras
cinco medidas permanecem confortavelmente acima do ponto de equilíbrio em toda a faixa de
âncoras testada. As cinco figuras abaixo mostram esse padrão em mais detalhe, da
comparação bruta de custo-benefício até como o próprio risco está distribuído pela
cidade.

\needspace{3.11in}
\begin{figure}[H]
\centering
\includegraphics[width=0.55\textwidth]{figures/brazil_companies_cost_benefit.png}
\caption{Resultados de custo-benefício, Companies}
\label{fig:companies-cb}
\end{figure}

A Figura \ref{fig:companies-cb} é a contraparte visual da Figura \ref{fig:residential-cb} e
o contraste é nítido: aqui os retângulos sobem bem acima da linha de equilíbrio, com o
armazenamento elevado de estoque chegando a 9,08 e a levee do distrito comercial a 6,44. Os
mesmos seis mecanismos, aplicados a uma base de exposição muito menor e mais concentrada,
produzem um portfólio que se paga várias vezes. Apenas o seguro fica abaixo da linha, e pelo
motivo estrutural discutido na Seção \ref{sec:sintese}, não por causa de como foi custeado.

\needspace{2.63in}
\begin{figure}[H]
\centering
\includegraphics[width=0.55\textwidth]{figures/brazil_companies_benefits_by_return_period.png}
\caption{Benefícios por período de retorno, Companies}
\label{fig:companies-rp}
\end{figure}

A Figura \ref{fig:companies-rp} repete o padrão de corte visto em Residential. A levee, com
corte no período de retorno de 20 anos, domina o benefício nos períodos de 5 e 10 anos e cai
a zero a partir de 20 anos. Daí em diante o seguro domina na camada de risco mais alta,
enquanto os mecanismos de transformação de MDD e PAA, o estoque elevado e o planejamento de
continuidade de negócios, contribuem uma faixa intermediária relativamente estável em todos
os períodos de retorno, já que nenhum deles tem um teto de proteção que o evento possa
exceder.

\needspace{3.32in}
\begin{figure}[H]
\centering
\includegraphics[width=0.48\textwidth]{figures/brazil_companies_waterfall_averted_risk.png}
\caption{Waterfall de risco evitado, Companies}
\label{fig:companies-wf}
\end{figure}

A Figura \ref{fig:companies-wf} completa o conjunto das três decomposições de risco. O risco
sobe de 14,3 milhões de dólares em 2025 para 49,6 milhões em 2050, com 9,1 milhões
atribuíveis ao crescimento econômico e 26,2 milhões à mudança climática, perto de 74\% clima
contra 26\% crescimento. Este é o terceiro ativo a cair na mesma faixa estreita apesar de
dados de exposição, calibração e medidas inteiramente independentes, o que constitui a
verificação de consistência interna mais forte disponível neste exercício.

\needspace{4.19in}
\begin{figure}[H]
\centering
\includegraphics[width=0.55\textwidth]{figures/brazil_companies_map_expected_annual_impact.png}
\caption{Mapa de risco anual esperado, Companies}
\label{fig:companies-eai}
\end{figure}

A Figura \ref{fig:companies-eai} mostra uma concentração dupla, em vez do agrupamento único
de Schools ou da dispersão difusa de Residential. Uma faixa densa corre ao longo do rio
pelos bairros Farrapos, Humaitá e Sarandi, coincidindo com o terreno mais baixo e mais
alagável, e um segundo núcleo, bem menor, fica no centro histórico de Porto Alegre, sua área
comercial e portuária. O fato de a atividade comercial se concentrar exatamente nos locais
mais expostos à inundação é o que dá a esse ativo sua alta razão risco-exposição e, por
consequência, o motivo de medidas direcionadas terem desempenho tão bom nele.

\needspace{4.99in}
\begin{figure}[H]
\centering
\includegraphics[width=0.55\textwidth]{figures/brazil_companies_map_rp100_depth_relocation.png}
\caption{Profundidade no período de retorno de 100 anos e pontos de relocação, Companies}
\label{fig:companies-depth}
\end{figure}

A Figura \ref{fig:companies-depth} marca as 20 empresas selecionadas para relocação contra a
mancha de profundidade de 100 anos. Quase todas caem dentro da faixa ribeirinha densa, com
apenas alguns pontos isolados fora dela, um caso intermediário entre o agrupamento compacto
de Schools e a ampla dispersão de Residential. Isso é coerente com a relocação devolver aqui
uma razão benefício-custo positiva mas estreita, de 1,32, o único resultado de Companies
sinalizado na Seção \ref{sec:sensibilidade} como sensível à escolha da âncora.

\subsection{Síntese comparativa entre os três ativos}
\label{sec:sintese}

O seguro se mantém consistentemente próximo de uma razão benefício-custo de 0,66 nos três
ativos. Trata-se de um reflexo mecânico do carregamento de custo de 1,5 vez aplicado
uniformemente, não de uma falha de modelagem; a razão depende só do carregamento de custo,
não da escala de risco. Vale declarar isso explicitamente a um tomador de decisão, porque um
gráfico de barras que coloca o seguro abaixo de todas as outras medidas convida à leitura de
que ele foi mal parametrizado, quando na verdade uma razão perto do inverso do carregamento
é exatamente o que uma camada de transferência de risco corretamente especificada deve
produzir.

Medidas de infraestrutura ampla, como levees, dredging e drenagem, apresentam razão
benefício-custo muito diferente conforme o ativo avaliado. Em Schools, essas medidas
protegem a cidade inteira, mas são avaliadas apenas contra a fatia estreita de escolas,
resultando em razões entre 0,25 e 1,60. Em Companies, a levee do distrito comercial,
avaliada contra a população que de fato protege, chega a 6,44. Em Residential, a mesma
lógica de base mais ampla se aplica à drenagem, mas sua escala de risco não deixa margem:
sua razão cai para 0,73. Ou seja, o efeito de avaliar contra uma base maior continua
elevando a razão benefício-custo em relação a Schools, mas não é suficiente, por si só,
para compensar um risco real muito menor do que o valor exposto contra o qual é medido.

A relocação não é uniformemente positiva entre os três ativos: em Schools a razão é 3,31;
em Companies é 1,32, positiva mas por uma margem estreita; em Residential é 0,48, não se
pagando. A abordagem conservadora adotada, relocando apenas 60\% do valor exposto nos
pontos de maior profundidade, ainda assim não é suficiente para justificar o custo de
relocação de Residential frente ao seu risco estimado. O padrão entre os três ativos
acompanha o quanto o risco de cada um é espacialmente concentrado, como mostram os mapas de
profundidade da Seção \ref{sec:resultados}: bem agrupado em Schools, intermediário em
Companies e difuso em Residential. A relocação recompensa a concentração, porque um número
fixo de pontos relocados captura uma fatia maior do risco total quando esse risco não está
espalhado de forma rala.

A proporção entre mudança climática e crescimento econômico no aumento do risco futuro é
notavelmente estável nos três ativos, entre 71\% e 75\% atribuível ao clima e entre 25\% e
29\% ao crescimento econômico. Esse resultado é um sinal de robustez do exercício: os três
ativos usam bases de exposição, calibração e medidas completamente diferentes, mas
compartilham o mesmo cenário de clima futuro e a mesma taxa de crescimento econômico de 2\%
ao ano. A estabilidade dessa proporção sugere que o cenário climático, e não escolhas
particulares de parametrização, é o principal motor do aumento do risco projetado.

Por fim, as medidas com corte de período de retorno, como drenagem e levee, sempre zeram
acima do seu teto de proteção. Esse padrão aparece com clareza nos gráficos de benefícios
por período de retorno de Residential e Companies: proteção total até o período de retorno
de 20 anos, contribuição nula depois. Trata-se de uma limitação conhecida desse mecanismo,
que oferece proteção completa ou nenhuma, sem gradação parcial, e vale a pena explicar com
clareza a um tomador de decisão antes que a medida seja adotada isoladamente.

\section{Análise de causa raiz: por que as medidas de Residential não se pagam}
\label{sec:causa-raiz}

O resultado de que nenhuma das seis medidas de Residential se paga, enquanto todas as seis
de Companies se pagam, motivou uma investigação dedicada, já que um resultado tão
surpreendente merece verificação antes de ser aceito como confiável. Duas verificações
foram feitas. Primeiro, um clone limpo do repositório upstream do Starter Pack confirmou
que não existe nenhum exemplo oficial de custo-benefício para Residential ou Companies em
lugar nenhum do upstream: apenas Schools traz um exemplo completo e calibrado. Não existe,
portanto, um resultado de Residential "que já funcionava antes" e que este exercício
teria feito regredir; esta é a primeira vez que esse cálculo é feito para esse ativo.
Segundo, a reexecução de notebooks descrita na Seção 2.2 confirmou que o próprio pipeline
reproduz seus números exatamente, descartando um bug de execução ou de ambiente.

Com a reprodutibilidade e a validade da linha de base confirmadas, a própria assimetria
ainda precisa de uma explicação, já que é grande o suficiente para que um leitor
razoavelmente suspeitasse de um erro de modelagem em vez de um achado genuíno. A análise a
seguir identifica duas causas que se somam, uma sobre como as seis medidas foram custeadas
e outra sobre a economia subjacente das duas bases de exposição, e quantifica cada uma em
vez de deixar a explicação qualitativa. A Tabela \ref{tab:causaraiz} apresenta a
comparação que motiva as duas: Residential e Companies diferem enormemente em valor
exposto, mas muito menos em risco climático total, e essa diferença por si só explica boa
parte do que se segue.

\begin{table}[H]
\centering
\caption{Comparação custo/risco, Residential versus Companies}
\label{tab:causaraiz}
\small
\begin{tabular}{lp{2.1cm}p{2.1cm}p{2.1cm}p{2.3cm}p{1.5cm}}
\toprule
Ativo & Valor exposto (2025) & Risco climático total & Risco / valor exposto & Soma dos 6 custos de medida & Custo / risco \\
\midrule
Residential & 69,312 bi USD & 470,267 M USD & 0,68\% & 1.402,760 M USD & 298\% \\
Companies & 3,656 bi USD & 828,925 M USD & 22,67\% & 898,569 M USD & 108\% \\
\bottomrule
\end{tabular}
\end{table}

\textbf{Causa 1: vários custos de medida de Residential são orçamentos de programa em
escala de cidade, grandes em relação ao seu risco.} Três das seis medidas de Residential,
Home flood-proofing (362,064 milhões de dólares), Community-based early warning (150,000
milhões de dólares) e Flood-resistant building codes (200,000 milhões de dólares), são
custeadas como programas em escala de cidade, pensados para alcançar toda a base
residencial, 69,3 bilhões de dólares de valor exposto espalhados por mais de 20.000 pontos
de grade. Apenas uma das seis medidas de Companies é custeada numa escala de programa
comparável. Essa assimetria de escala de custo, independente de qualquer questão de
calibração de curva, já explica boa parte da diferença de resultado benefício-custo entre
os dois ativos.

\textbf{Causa 2: o descasamento de escala também é estrutural.} Residential tem cerca de 19
vezes mais valor exposto que Companies, mas seu risco climático total é apenas
marginalmente maior em termos absolutos. Isso torna a intensidade de risco de Residential,
risco como fração do valor exposto, cerca de 33 vezes menor que a de Companies (0,68\%
contra 22,67\%). A soma dos seis custos de medida de Residential, 1,4 bilhão de dólares,
quase três vezes seu próprio risco total, simplesmente não cabe dentro do risco que
endereça, mesmo onde nenhum custo individual é incomumente alto em termos absolutos.

A diferença em quanto a curva de literatura sem ajuste superestimou a perda observada em
2024, cerca de 20,7 vezes para Residential contra 2,8 vezes para Companies (Seção
\ref{sec:curva-profundidade}), remonta à própria curva JRC: a tabela "South America
residential" se encaixou muito pior na perda real de Porto Alegre em 2024 do que a tabela
"South America commerce". Não é possível afirmar com confiança, a partir dos dados
atualmente disponíveis, se isso se deve à curva ser um ajuste ruim para o parque
residencial de Porto Alegre especificamente, ou a esse parque ser em si atípico em relação
à média regional que a curva representa; resolver isso exigiria um mapa de profundidade do
evento de 2024, a mesma limitação já apontada na Seção \ref{sec:curva-profundidade}.

Este não é um erro computacional. É um achado legítimo, ainda que desconfortável: na
melhor estimativa disponível de risco real, a adaptação de Residential, do jeito em que foi
delineada e custeada neste exercício, não é economicamente justificada, enquanto o
portfólio de Companies continua justificado, com margens reduzidas. A implicação prática não
é que a adaptação residencial deva ser abandonada, e sim que ela precisa ser redimensionada:
as medidas aqui foram especificadas em escala de cidade, e a análise sugere que o argumento
a favor delas precisa ser feito ou numa escala menor e mais bem direcionada, ou em bases
que não o dano econômico direto evitado, como deslocamento, saúde ou continuidade de
serviços, nenhuma das quais este modelo precifica.

\section{Discussão e conclusões}

O exercício confirma que estender um exemplo de custo-benefício de um ativo para outros
dentro do mesmo estudo de caso não é uma tarefa mecânica de repetição de código. A
calibração de um novo ativo exige um espaço de busca próprio, sensível à escala de
exposição daquele ativo, e a integração entre calibração e análise de custo-benefício
exige atenção explícita à unidade física do hazard usado em cada etapa: Residential e
Companies são calibrados contra um hazard de duração para reconstrução de evento, mas a
análise de custo-benefício roda sobre um hazard de profundidade, então usa uma curva
separada, baseada em profundidade (Seção \ref{sec:curva-profundidade}). Uma unidade física
declarada compatível com o hazard efetivamente usado é necessária, mas não suficiente: a
amplitude da curva de literatura, testada sem ajuste contra a perda real observada em 2024,
superestimava essa perda em 20,7 vezes para Residential e 2,8 vezes para Companies, motivo
pelo qual sua amplitude é calibrada localmente em vez de aceita ao pé da letra. A lição mais
geral é que a unidade física de uma curva de dano e sua magnitude são duas propriedades
distintas a verificar; uma curva pode ser dimensionalmente correta e ainda assim produzir
resultados substancialmente incorretos, silenciosamente, porque o resultado parece
plausível o suficiente para não levantar suspeita, a menos que se confronte deliberadamente
com uma perda real observada.

Do ponto de vista de decisão, o quadro final é mais heterogêneo do que uma conclusão única
do tipo "as medidas de adaptação valem a pena". Para Companies, o portfólio de seis
medidas faz sentido econômico, com margens modestas mas todas acima do ponto de equilíbrio.
Para Residential, nenhuma das seis medidas
parametrizadas se paga ao risco real estimado: nem a infraestrutura ampla (drenagem), nem
a relocação parcial, nem as medidas de MDD ou PAA. Isso não significa que a adaptação não
valha a pena para o setor residencial de Porto Alegre; significa que as medidas, do jeito
em que foram parametrizadas neste exercício, nesses custos e nessa escala, não se
justificam frente à melhor estimativa disponível de risco. Um refinamento futuro
precisaria repensar a escala das intervenções, por exemplo cobrindo só os pontos de maior
risco em vez da cidade inteira, ou buscar reduções de custo, não simplesmente aceitar os
números correntes. O seguro paramétrico, nos três ativos, continua sendo estruturalmente o
mecanismo de razão benefício-custo mais baixa (aproximadamente 0,66), um custo de
transferência de volatilidade que vale a pena explicar a um tomador de decisão como algo
esperado, não como uma medida malparametrizada.

A limitação mais importante deste exercício é que a curva de dano final, embora ancorada
na perda real de 2024, foi calibrada contra um único ponto de dado agregado, a perda total
do evento, não contra uma distribuição espacial de perdas como a que orientou a
calibração do hazard de duração na Seção \ref{sec:calibration}. O período de retorno 200 usado como proxy da
profundidade de 2024 foi escolhido por semelhança de área alagada entre dois modelos de
hazard distintos, não por uma correspondência física direta; é a melhor aproximação
disponível, não uma calibração com o mesmo rigor da Seção \ref{sec:calibration}. Qualquer refinamento futuro
do estudo de caso do Brasil deveria priorizar a obtenção ou reconstrução de um mapa de
profundidade do evento de 2024, para permitir uma calibração direta e espacialmente
distribuída, da mesma forma que já existe para o hazard de duração.

\section{Recomendações para replicação}

Todo número neste relatório é reproduzível a partir dos dados e do código entregues junto
com ele, sem precisar reconstruir nada do zero. Os cinco notebooks que produzem os
resultados da Seção \ref{sec:resultados}, a calibração em \texttt{06b} e as quatro
execuções de custo-benefício em \texttt{09} a \texttt{12}, podem cada um ser reexecutado
de forma independente uma vez que o ambiente descrito na Seção 2.1 esteja configurado, e
cada um escreve suas próprias figuras e totais impressos no lugar, então uma discrepância
em qualquer número pode ser rastreada até a célula específica do notebook que a produziu.
Os comandos abaixo assumem um terminal aberto na raiz deste repositório, com o ambiente
\texttt{climada\_env} já criado e ativado.

\begin{verbatim}
conda activate climada_env
cd notebooks/
jupyter nbconvert --to notebook --execute --inplace \
    06b_calibration_residential_companies.ipynb
jupyter nbconvert --to notebook --execute --inplace 09_costbenefit.ipynb
jupyter nbconvert --to notebook --execute --inplace \
    10_costbenefit_additional_measures.ipynb
jupyter nbconvert --to notebook --execute --inplace \
    11_costbenefit_residential.ipynb
jupyter nbconvert --to notebook --execute --inplace \
    12_costbenefit_companies.ipynb
\end{verbatim}

Todos os arquivos de entrada necessários, incluindo os Entity files calibrados e com
medidas e os arquivos HDF5 de relocação, já estão disponíveis em \texttt{data/}. Não é
preciso rodar os scripts \texttt{.py} numerados para reproduzir os resultados deste
relatório; eles servem apenas como registro de como cada arquivo foi construído e estão
documentados em \texttt{notebooks/README.md}. Rodá-los só é necessário para reconstruir os
insumos do zero ou adaptar parte do exercício.

Para reproduzir ou estender esta análise dentro do estudo de caso do Brasil, cinco passos
gerais se aplicam. Primeiro, calibrar o ativo de interesse seguindo o método dos notebooks
\texttt{06a} e \texttt{06b}, ajustando o espaço de busca à escala de exposição daquele
ativo, pois o espaço de busca original não é transferível entre ativos de escalas
diferentes. Segundo, e mais importante, confirmar que a unidade física do hazard usado na
calibração é a mesma unidade física do hazard que será consumido na análise de
custo-benefício; esse deveria ser um passo obrigatório antes de qualquer calibração, já que
um descasamento de unidade não gera exceções nem mensagens de erro, apenas resultados
numericamente plausíveis e substancialmente incorretos. Terceiro, se uma curva de
literatura for adotada para resolver um descasamento de unidade ou domínio, validar sua
magnitude contra o melhor dado local disponível antes de confiar em valores absolutos,
mesmo que seja uma única observação agregada, como faz a Seção
\ref{sec:curva-profundidade}; não assumir que uma curva dimensionalmente correta já está
calibrada. Quarto, desenhar as medidas de adaptação de
modo a cobrir mecanismos distintos do CLIMADA, já que variar o mecanismo, e não apenas o
valor numérico dos parâmetros, é o que torna a análise informativa para o tomador de
decisão. Quinto, refinar a parametrização segundo os três critérios de aceite da Seção
\ref{sec:refinamento} antes de considerar os números finais, e resistir à tentação de
reajustar os custos uma segunda vez só para fazer um resultado desconfortável parecer
melhor.

\bibliographystyle{plainnat}
\bibliography{references}

\appendix

\section{Notas técnicas e recomendações para o Starter Pack}
\label{app:notas}

As notas a seguir foram compiladas durante a extensão do pack para este exercício, para
discussão com a equipe da UNU-EHS e para outros participantes trabalhando no estudo de
caso do Brasil ou estendendo o seu próprio. Elas cobrem pontos de fricção na estrutura de
dados, nos utilitários de plotagem e no fluxo de calibração do Starter Pack que não eram
óbvios a partir do material de ensino existente, mas que afetaram materialmente quanto
tempo este exercício levou e quão fácil é confiar nos resultados. Algumas são correções
pontuais de uma linha; outras, como a distinção entre uma curva dimensionalmente correta e
uma curva localmente validada, apontam para um hábito metodológico mais amplo que vale a
pena incorporar à documentação do pack, em vez de deixar que cada participante o redescubra
de forma independente. Elas estão ordenadas aproximadamente pela ordem em que um
participante típico estendendo o pack provavelmente as encontraria.

\begin{enumerate}
\item \textbf{Medidas de troca de Exposures interagem mal com crescimento econômico, e o
pack não avisa sobre isso.} A coluna \texttt{assets file} da medida aponta para um único
arquivo HDF5, mas \texttt{CostBenefit.calc} aplica as medidas do Entity presente ao
cenário presente e as do Entity futuro ao cenário futuro. Se ambos os Entities apontarem
para o mesmo arquivo HDF5 de valores presentes, o benefício da medida inclui silenciosamente
a remoção de décadas de crescimento econômico do cenário futuro adaptado, superestimando-o
substancialmente. Este exercício resolveu isso escrevendo um segundo arquivo HDF5, já
escalado pelo crescimento, e reapontando a medida do Entity futuro em código. Um aviso
curto no notebook de custo-benefício, ou um exemplo trabalhado, pouparia outros
participantes dessa armadilha.

\item \textbf{Nenhum exemplo trabalhado deriva os parâmetros de seguro a partir do próprio
modelo.} O material de treinamento introduz \texttt{Impact.calc\_freq\_curve()} e os
parâmetros de transferência de risco (attach e cover) separadamente, mas não os conecta.
Este exercício definiu attach e cover a partir da curva de excedência presente e não
adaptada, o que pareceu o padrão natural a ensinar; vale a pena incluir isso no notebook,
já que constantes arbitrárias de primeira estimativa produzem prêmios e benefícios muito
diferentes.

\item \textbf{Caminhos relativos na coluna \texttt{assets file} resolvem contra o diretório
de trabalho do processo, não contra a localização do próprio Entity file.} Isso funciona
enquanto os notebooks rodam a partir da pasta \texttt{notebooks/}, mas quebra assim que um
Entity file é carregado de qualquer outro lugar. Essa fricção se repetiu duas vezes durante
as próprias correções desta sessão: uma ao copiar um Entity file corrigido entre duas
estruturas de pasta com nomes diferentes, e outra ao renomear essas mesmas pastas para
publicação pública, ambas quebrando a medida de relocação silenciosamente até serem
reexecutadas e conferidas. Resolver \texttt{assets file} relativo à localização do próprio
Entity file, em vez do diretório de trabalho do processo, eliminaria essa classe inteira
de bug.

\item \textbf{Uma curva dimensionalmente correta não é o mesmo que uma curva validada.} O
pack é cuidadoso com a incerteza de calibração onde calibra explicitamente, mas a curva
baseada em profundidade que vem com o Entity de Schools chega pronta, sem nota sobre sua
procedência nem sobre se sua amplitude alguma vez foi conferida contra dados locais. Adotar
a curva de literatura equivalente para Residential e Companies e depois conferi-la contra a
perda real de 2024 em Porto Alegre (Seção \ref{sec:curva-profundidade}) mostrou a versão
sem ajuste superestimando essa perda em cerca de 20,7 vezes para Residential e 2,8 vezes
para Companies. Rotular a curva que vem no pack com sua fonte e acrescentar uma única frase
ao lado de qualquer sugestão de curva de literatura, afirmando que ela é um padrão regional
cuja amplitude deve ser validada contra dados locais disponíveis antes de confiar em
valores absolutos, ajudaria outros participantes a capturar isso bem mais cedo.

\item \textbf{Os limites de julgamento especializado que resolvem a degenerescência da
calibração são específicos do ativo, e nada avisa quando eles deixam de se aplicar.} O pack
trata isso bem: o \texttt{06a} explica claramente que ajustar dois parâmetros a uma única
observação deixa toda uma família de soluções igualmente boas, e resolve isso com
julgamento especializado declarado para Schools, \texttt{max\_damage > 0,7} porque várias
escolas foram inteiramente reconstruídas, e \texttt{max\_intensity < 30} lido no gráfico do
espaço de parâmetros como consequência dessa primeira escolha. A dificuldade é que o
segundo limite deriva do primeiro, e apenas o segundo sobrevive à transferência para outro
ativo: Residential e Companies calibram para razões de dano de 0,028 e 0,294, então
\texttt{max\_damage > 0,7} é impossível por construção, enquanto \texttt{max\_intensity <
30} continua rodando e continua selecionando um resultado silenciosamente. Um participante
adaptando o \texttt{06a} para um novo ativo herda, portanto, um limite numérico despido de
sua justificativa. Reescrever a célula de julgamento especializado como uma checklist
curta, perguntando quais limites são específicos do ativo e que evidência os definiria para
o novo ativo, tornaria essa dependência visível no ponto em que ela importa.

\item \textbf{A barra de cores automática de \texttt{plot\_basemap\_eai\_exposure()}
colide com um título customizado quando os valores excedem $10^6$.} O matplotlib adiciona
um rótulo de notação científica acima da barra de cores, na mesma altura de um título
centralizado; para Schools, valores abaixo de 200.000, isso nunca aparece, mas Residential
e Companies, valores de até vários milhões de dólares, disparam o problema. A correção é
capturar o eixo da barra de cores e forçar um \texttt{ScalarFormatter} simples, em vez de
brigar com o comprimento ou posição do título. Também, chamar
\texttt{ax.set\_title(..., loc="left")} depois de uma função que já define seu próprio
título centralizado cria um segundo título independente em vez de substituir o primeiro,
já que o matplotlib guarda os títulos esquerdo, central e direito separadamente; chamar
\texttt{set\_title} sem o argumento \texttt{loc} substitui o título existente.

\item \textbf{Falha de kernel específica do Windows, não reproduzida no Linux.} No
Windows, o build do \texttt{pyzmq}/\texttt{libzmq} do conda-forge pode fazer todo
\texttt{jupyter nbconvert --execute} falhar com \texttt{Bad file descriptor}
(\texttt{epoll.cpp:73}); substituir o \texttt{pyzmq} pelo wheel do PyPI
(\texttt{pip install --force-reinstall pyzmq==26.2.0}) restaura a comunicação com o
kernel. Os próprios notebooks deste exercício foram reexecutados e verificados em Linux
sem encontrar esse problema, então ele é específico do build Windows do pacote
conda-forge, não dos notebooks em si; ainda vale uma linha no FAQ de instalação para
participantes que trabalham no Windows.

\item \textbf{A razão benefício-custo do seguro é aproximadamente um sobre o fator de
custo por construção, e vale a pena dizer isso explicitamente.} Com um carregamento de
prêmio de 1,5, a razão fica perto de 0,66 independente da camada. O valor da medida está
na transferência de volatilidade, proteção contra anos raros que quebram o orçamento, não
economia de valor esperado; uma frase no material de ensino evitaria a leitura equivocada
de que "o seguro parece ruim" no gráfico de barras.

\item \textbf{A etapa de exportação Excel do \texttt{06a} é manual e rotula errado a
unidade de intensidade.} A célula final de \texttt{06a\_calibration\_schools.ipynb}
imprime as funções calibradas com a coluna \texttt{intensity\_m} e
\texttt{intensity\_unit: "meter"}, embora a calibração tenha rodado sobre duração de
inundação em dias, e a tabela precisa então ser colada manualmente na planilha calibrada.
O notebook \texttt{06b} automatizou essa exportação com openpyxl e escreveu a unidade
como \texttt{days}. Também vale esclarecer que o Entity de Schools calibrado que vem no
pack usa uma curva baseada em profundidade (metros), consistente com os rasters de período
de retorno do JRC, que não é a curva baseada em duração que o próprio \texttt{06a}
produz; a conexão entre as duas mereceria ser explicada explicitamente no material de
ensino.
\end{enumerate}

\end{document}
